\documentclass[11pt,letterpaper]{article}

% ============================================================================
% PACKAGES
% ============================================================================
\usepackage[margin=1in, top=1.5in, headheight=85pt]{geometry}
\usepackage{graphicx}
\usepackage{fancyhdr}
\usepackage{lastpage}
\usepackage{enumitem}
\usepackage{xcolor}
\usepackage{hyperref}
\usepackage{tabularx}

% ============================================================================
% DOCUMENT VARIABLES
% ============================================================================
\newcommand{\UniqueID}{DM-2026-003}
\newcommand{\DocumentDate}{January 21, 2026}
\newcommand{\AuthorName}{SendCUIEmail Development Team}
\newcommand{\AuthorTitle}{Software Engineering}
\newcommand{\ToField}{Project Record}
\newcommand{\SubjectField}{Password/Key Transmission Requirements per NIST Guidelines}
\newcommand{\OPRField}{SendCUIEmail Project}

% ============================================================================
% DOCUMENT CONFIGURATION
% ============================================================================
\hypersetup{
    colorlinks=true,
    linkcolor=blue,
    urlcolor=blue,
    pdftitle={DM-2026-003 Password Transmission},
    pdfauthor={SendCUIEmail Project},
}

\setlength{\parindent}{0pt}
\setlength{\parskip}{6pt}

\setlist[enumerate]{leftmargin=0.5in, labelwidth=0.25in, labelsep=0.1in, align=left, itemsep=12pt, topsep=12pt}
\setlist[enumerate,1]{label=\arabic*.}
\setlist[itemize]{leftmargin=0.5in, itemsep=3pt, topsep=6pt}

% ============================================================================
% HEADER AND FOOTER
% ============================================================================
\pagestyle{fancy}
\fancyhf{}
\fancyhead[L]{\textbf{SendCUIEmail}}
\fancyhead[R]{\parbox[b]{3.5in}{\raggedleft\large\textbf{Decision Memorandum}}}
\renewcommand{\headrulewidth}{0pt}
\fancyfoot[L]{\small \UniqueID}
\fancyfoot[C]{\small Page \thepage\ of \pageref{LastPage}}
\fancyfoot[R]{\small \DocumentDate}
\renewcommand{\footrulewidth}{0.4pt}

\newcommand{\dmsection}[1]{\item \textbf{#1}\par}

% ============================================================================
% BEGIN DOCUMENT
% ============================================================================
\begin{document}

\hfill \UniqueID

\vspace{0.25in}

\underline{\hspace{2.5in}}\\[2pt]
\AuthorName\\
\AuthorTitle

\vspace{0.25in}

\textbf{DATE:} \DocumentDate

\textbf{TO:} \ToField

\textbf{SUBJECT:} \SubjectField

\textbf{OFFICE OF PRIMARY RESPONSIBILITY:} \OPRField

\vspace{0.25in}

\begin{enumerate}

% ----------------------------------------------------------------------------
\dmsection{Purpose}

This memorandum documents NIST requirements for encryption key and password transmission, explains why SendCUIEmail exists, and provides guidance on when certificate-based encryption should be used instead.

% ----------------------------------------------------------------------------
\dmsection{Background}

SendCUIEmail provides password-based file encryption for situations where certificate-based encryption is not feasible. This document clarifies when each approach is appropriate and the security implications of password transmission methods.

% ----------------------------------------------------------------------------
\dmsection{NIST Guidance on Out-of-Band Authentication}

\textbf{NIST SP 800-63B Section 5.1.3.1} explicitly states:

\begin{quote}
``Methods that do not prove possession of a specific device, such as voice-over-IP (VOIP) or email, \textbf{SHALL NOT} be used for out-of-band authentication.''
\end{quote}

An out-of-band authenticator is defined as:
\begin{quote}
``A physical device that is uniquely addressable and can communicate securely with the verifier over an \textbf{independent communications channel}.''
\end{quote}

\textbf{Acceptable Out-of-Band Channels:}
\begin{itemize}
    \item SMS to mobile phone (proves possession of SIM card)
    \item Voice call over PSTN (Public Switched Telephone Network)
    \item In-person delivery
    \item Encrypted instant messaging (if it proves device possession)
\end{itemize}

\textbf{NOT Acceptable Out-of-Band Channels:}
\begin{itemize}
    \item Email (does not prove device possession)
    \item VoIP (does not prove device possession)
    \item Same communication channel as encrypted data
\end{itemize}

% ----------------------------------------------------------------------------
\dmsection{Security Hierarchy for CUI Email Transmission}

The following table presents the security options in order of preference:

\begin{tabularx}{\textwidth}{|c|l|X|}
\hline
\textbf{Rank} & \textbf{Method} & \textbf{Description} \\
\hline
1 & S/MIME (PIV/CAC to PIV/CAC) & Both parties have smart cards. Email signed with sender's private key, encrypted with recipient's public certificate. \textbf{Gold standard.} \\
\hline
2 & Hybrid (SendCUIEmail + S/MIME password) & Files: AES-256 password-encrypted. Password: S/MIME encrypted email. Compliant when recipient has PIV/CAC. \\
\hline
3 & SendCUIEmail + Out-of-band password & Files: AES-256 password-encrypted. Password: Phone/SMS/in-person. Compliant per NIST 800-63B. \\
\hline
4 & SendCUIEmail + Unencrypted email password & \textbf{NOT COMPLIANT} per NIST 800-63B. Convenience feature with warnings. \\
\hline
\end{tabularx}

\vspace{12pt}

\textbf{Key Insight:} SendCUIEmail exists for scenarios where S/MIME is not available. When both parties have PIV/CAC, native S/MIME is preferred.

% ----------------------------------------------------------------------------
\dmsection{Option 1: Native S/MIME (PIV/CAC to PIV/CAC)}

When \textbf{both sender and recipient} have PIV/CAC cards, native S/MIME provides the highest security with no password transmission required:

\textbf{Security Properties:}
\begin{itemize}
    \item \textbf{Confidentiality}: Encrypted with recipient's public certificate---only their PIV/CAC can decrypt
    \item \textbf{Authenticity}: Signed with sender's private key---proves who sent it
    \item \textbf{Non-repudiation}: Sender cannot deny sending (signature tied to identity)
    \item \textbf{Integrity}: Any tampering breaks the cryptographic signature
\end{itemize}

\textbf{Implementation (Outlook):}
\begin{enumerate}
    \item Compose email with CUI files attached
    \item Click \textbf{Sign} (uses your PIV/CAC private key)
    \item Click \textbf{Encrypt} (uses recipient's public certificate)
    \item Send
\end{enumerate}

\textbf{When to Use:}
\begin{itemize}
    \item Both parties have PIV/CAC with valid certificates
    \item Recipient's certificate is accessible (same organization, directory service, or previously exchanged)
    \item Regular communication with established recipients
\end{itemize}

\textbf{Note:} In this scenario, SendCUIEmail is \textbf{not needed}---use Outlook's native S/MIME instead.

% ----------------------------------------------------------------------------
\dmsection{Option 2: Why SendCUIEmail Uses Password-Based Encryption}

SendCUIEmail uses password-based (symmetric) encryption rather than certificate-based encryption because it is designed for recipients who \textbf{do not have} PIV/CAC or whose certificates are not accessible.

\textbf{Why Out-of-Band Requirements Don't Apply:}

The NIST SP 800-63B out-of-band requirement applies to \textbf{shared secrets} (passwords, symmetric keys). Certificate-based encryption uses \textbf{asymmetric cryptography}---no shared secret exists:

\begin{tabularx}{\textwidth}{|l|l|X|}
\hline
\textbf{Encryption Type} & \textbf{What's Transmitted} & \textbf{Out-of-Band Needed?} \\
\hline
Password-based & Shared secret (password) & Yes---per NIST 800-63B \\
\hline
Certificate-based & Public key only & No---public keys aren't secret \\
\hline
\end{tabularx}

\vspace{12pt}

\textbf{How Certificate-Based Encryption Works:}
\begin{itemize}
    \item Sender encrypts using recipient's \textbf{public certificate} (can be sent via email---not secret)
    \item Only recipient can decrypt using their \textbf{private key} (never leaves PIV/CAC card)
    \item \textbf{No shared secret} exists, so no out-of-band transmission is required
    \item Provides non-repudiation and authentication
\end{itemize}

\textbf{Standards:}
\begin{itemize}
    \item \textbf{S/MIME} (Secure/Multipurpose Internet Mail Extensions)
    \item \textbf{OpenPGP} (Pretty Good Privacy)
    \item Both are referenced in NIST SP 800-45 for email security
\end{itemize}

\textbf{When to Use Certificate-Based Encryption:}
\begin{itemize}
    \item Recipient has PIV/CAC card with valid certificate
    \item Certificate can be obtained from directory service or key server
    \item Organization has PKI infrastructure in place
    \item Regular communication with same recipients
\end{itemize}

% ----------------------------------------------------------------------------
\dmsection{When SendCUIEmail Is Appropriate}

SendCUIEmail's password-based encryption is designed for situations where certificate-based encryption is \textbf{not feasible}:

\begin{itemize}
    \item Recipient does not have PIV/CAC card
    \item Recipient's certificate is not accessible (different organization, no directory access)
    \item One-time or infrequent communication with external parties
    \item Recipient's email client does not support S/MIME
    \item Emergency situations where certificate exchange cannot be completed in time
\end{itemize}

% ----------------------------------------------------------------------------
\dmsection{Hybrid Approach: S/MIME Encrypted Password Email}

A compliant alternative uses \textbf{both} encryption methods:

\begin{enumerate}
    \item \textbf{Email \#1}: Files encrypted with SendCUIEmail (password-based AES-256)
    \item \textbf{Email \#2}: Password in message body, \textbf{email itself S/MIME encrypted} with recipient's PIV/CAC certificate
\end{enumerate}

\textbf{Why This Is Compliant:}
\begin{itemize}
    \item The password (human-readable/enterable) is protected by asymmetric encryption
    \item \textbf{Encrypted} with recipient's public certificate---only they can decrypt
    \item \textbf{Signed} with sender's private key (from sender's PIV/CAC)---proves authenticity
    \item Same channel (email) is acceptable when the secret is cryptographically protected
    \item The encrypted files in Email \#1 can be forwarded, stored, or used on systems without S/MIME
\end{itemize}

\textbf{When to Use This Approach:}
\begin{itemize}
    \item Recipient has PIV/CAC but may need to forward/store files on non-S/MIME systems
    \item Need audit trail of password delivery (S/MIME provides non-repudiation)
    \item Organization policy requires all CUI via email be S/MIME encrypted
\end{itemize}

\textbf{Outlook Implementation:} When composing the password email, click ``Sign'' and ``Encrypt'' in Outlook to S/MIME sign and encrypt before sending.

\textbf{Future Enhancement:} PIV/CAC certificate access is already implemented in the PDFSigner project. This functionality could be integrated into SendCUIEmail to automate S/MIME signing/encryption of the password email when PIV/CAC is available.

% ----------------------------------------------------------------------------
\dmsection{Password Transmission Without S/MIME}

When \textbf{not} using S/MIME encryption for the password email, the password \textbf{MUST} be transmitted via an out-of-band channel per NIST SP 800-63B:

\begin{tabularx}{\textwidth}{|l|l|X|}
\hline
\textbf{Method} & \textbf{Compliant?} & \textbf{Notes} \\
\hline
Phone call (PSTN) & Yes & Proves possession of phone number \\
\hline
SMS text message & Yes & Proves possession of SIM card \\
\hline
In-person & Yes & Direct verification of identity \\
\hline
Separate email & \textbf{No} & Does not prove device possession \\
\hline
Same email & \textbf{No} & No separation at all \\
\hline
\end{tabularx}

\vspace{12pt}

\textbf{Note:} SendCUIEmail generates a \texttt{Password\_Email\_*.msg} draft as a \textbf{convenience feature} for operational flexibility. Users should understand this does \textbf{not} meet NIST out-of-band requirements. The tool displays warnings to this effect.

% ----------------------------------------------------------------------------
\dmsection{Implementation in SendCUIEmail}

The tool implements the following to guide users toward compliant behavior:

\begin{itemize}
    \item \textbf{Primary guidance}: ``Share password via phone, text, or in-person''
    \item \textbf{Secondary option}: Password email draft with security warning
    \item \textbf{Warning in password email}: Explicitly states separate channel is preferred
    \item \textbf{No auto-send}: All email drafts require manual recipient entry and review
\end{itemize}

% ----------------------------------------------------------------------------
\dmsection{Decision}

\begin{enumerate}
    \item \textbf{Certificate-based encryption} (S/MIME with PIV/CAC) is the preferred method when feasible
    \item When using password-based encryption, passwords \textbf{SHALL} be transmitted via out-of-band channel (phone, SMS, in-person)
    \item The password email feature is retained for operational flexibility but includes prominent warnings about NIST non-compliance
    \item Documentation clearly states these requirements
\end{enumerate}

% ----------------------------------------------------------------------------
\dmsection{References}

\begin{itemize}
    \item NIST SP 800-63B: Digital Identity Guidelines - Authentication and Lifecycle Management
    \item NIST SP 800-57 Part 1: Recommendation for Key Management
    \item NIST SP 800-45: Guidelines on Electronic Mail Security
    \item NIST SP 800-171: Protecting Controlled Unclassified Information
    \item FIPS 201: Personal Identity Verification (PIV) of Federal Employees and Contractors
\end{itemize}

\end{enumerate}

\end{document}
